\section{Limitations and Conclusion}\label{sec:conclusion}

\paragraph{Limitations} The model makes several simplifications.
\begin{itemize}
\item \emph{One secret interrupt.} The NOP handler exists only as the disk handler's
  indistinguishable partner. Several secret interrupts would need a richer
  indistinguishability argument.
\item \emph{Equal handler lengths.} Every handler runs for the same number of steps,
  which is what lets every mask bit be public. Variable-length handlers would force
  the masks private, and the proof of \m{fv_NI} would have to reason across diverging
  branches. A real system would reach a fixed length by padding.
\item \emph{Cooperative processes, no nesting.} A selected process takes exactly one
  step; preemption within a step is not modelled. Handlers cannot be interrupted,
  since all masks are set while one runs.
\item \emph{Abstract time.} One output step is one unit of time, for every process
  and handler alike. Caches, memory contention and other microarchitectural channels
  are outside the model.
\item \emph{One-step handlers.} At $\mathit{runtime} = 1$ a slice loses its last
  handler run (\cref{sec:restore}); the result still holds, but the policy behaves
  less usefully there.
\end{itemize}

\paragraph{Conclusion} A secret interrupt can leak through the schedule alone, and
servicing interrupts only in publicly timed, padded slices closes the channel. We
proved both facts in Rocq, the second for any non-interfering scheduler and user
processes, by composing generic non-interference theorems over the structure of the
model. The proof did more than check the policy: it fixed the classification of the
interrupt controller bit by bit, separated hiding what a handler computed from hiding
that it ran, and exposed a proof device that was quietly carrying behaviour, which we
moved into the policy and then proved inert.

% TODO: future work, e.g. several secret interrupts, variable-length handlers with
% padding, or connecting the model to a real interrupt controller.
